\subsection{کالبدشکافی ماژول \lr{gated\_linear\_attention}}

\subsubsection*{۱. فلسفه طراحی، پارادایم معماری و تمثیل ملموس}
مکانیزم نوآورانه \lr{gated\_linear\_attention} پاسخی بنیادین به یکی از حیاتی‌ترین چالش‌های مهندسی مدل‌های هوش مصنوعی مدرن است: \textbf{چگونه می‌توان ظرفیت بازنمایی زبانی مدل را در بالاترین سطح نگه داشت و هم‌زمان هزینه‌های سنگین پردازشی و بن‌بست‌های حافظه سخت‌افزار را مهار کرد؟}

\textbf{تمثیل ملموس و مثال کاربردی:}
برای فهم عمیق عملکرد \lr{gated\_linear\_attention} \lr{(GLA)}، یک مشاور ارشد با پتانسیومتر (دیمر) هوشمند تنظیم تمرکز را در نظر بگیرید. در توجه استاندارد، تمام مکالمات گذشته با وضوحی یکسان ثبت می‌شوند که سبب سرریز ذهن می‌گردد؛ اما در \lr{GLA}، برای هر بعد پنهان یک دیمر اختصاصی وابسته به محتوا تعبیه شده است. هرگاه موضوع گفتگو تغییر کند، دیمر ابعاد نامربوط را کم‌نور کرده و به سرعت وضعیت پنهان را برای دریافت کلیدواژه‌های تازه پاکسازی می‌کند، در حالی که نکات راهبردی پایدار با ضریب بالا در حافظه بازگشتی باقی می‌مانند.

\subsubsection*{۲. فرمول‌بندی و استنتاج عمیق ریاضی \lr{(Mathematical Derivation)}}
در این مکانیزم، توجه خطی با ماتریس قطری زوال داده‌محور فرمول‌بندی می‌گردد:
\begin{align}
\mathbf{S}_t &= \operatorname{diag}(\boldsymbol{\alpha}_t) \mathbf{S}_{t-1} + \mathbf{k}_t \mathbf{v}_t^T \\
\mathbf{o}_t &= (\mathbf{q}_t \mathbf{S}_t) \odot \mathbf{g}_t
\end{align}
که در آن بردار گیت فراموشی $\boldsymbol{\alpha}_t = \exp(-\operatorname{softplus}(\mathbf{W}_\alpha \mathbf{x}_t))$ پایداری لیاپانوف سیستم بازگشتی را تضمین کرده و گیت خروجی $\mathbf{g}_t = \sigma(\mathbf{W}_g \mathbf{x}_t)$ شدت انتقال تانسور حاصله به لایه بعد را تنظیم می‌کند.

\subsubsection*{۳. کالبدشکافی الگوریتمی و جریان داده \lr{(Algorithmic Tensor Trace)}}
در پیاده‌سازی سورس پایتون در فایل ماژول \lr{gated\_linear\_attention.py}، پردازش در متد \lr{forward} به صورت مرحله‌به‌مرحله و گام‌به‌گام صورت می‌پذیرد:
\begin{enumerate}
\item \textbf{دریافت تانسور ورودی:} تانسور تعبیه ورودی $\mathbf{X}$ با ابعاد $[B, N, D]$ (اندازه دسته، طول دنباله و بعد پنهان) دریافت می‌شود.
\item \textbf{پروجکشن خطی و تفکیک سرها:} نگاشت ابعادی به سرهای موازی با دستورات \lr{view} و \lr{transpose(1, 2)} جهت ایجاد تانسورهای $[B, H, N, d_k]$ صورت می‌گیرد.
\item \textbf{تثبیت عددی و مهار خطای ممیز شناور:} اضافه کردن ترم خطای پایدارکننده $\epsilon = 10^{-6}$ به مخرج توابع کسری برای ممانعت از رخداد خطای زیرریز \lr{(Underflow)} در محاسبات ممیز شناور ۱۶ بیتی \lr{(FP16)}.
\item \textbf{کاهش ابعاد و تجمیع اطلاعات:} اجرای عملگر تخصصی معماری با حذف درایه‌های زائد و استخراج مؤلفه‌های اصلی اطلاعات.
\item \textbf{بازگشت به بعد اصلی و پروجکشن پایانی:} تانسور به شکل $[B, N, D]$ بازآرایی شده و با یک لایه خطی پایانی به لایه پسین تزریق می‌گردد.
\end{enumerate}

\subsubsection*{۴. تحلیل همجوشی سخت‌افزاری و ترافیک حافظه \lr{(Hardware \& Roofline Analysis)}}
از دیدگاه مدل سقف کارایی \lr{(Roofline Model)} و معماری پردازنده‌های گرافیکی مدرن (نظیر کارت گرافیک \lr{NVIDIA RTX 4070}):
\begin{itemize}
\item \textbf{مهار عطش پهنای باند حافظه } \lr{(Memory-Bound Relief)}\rl{:}  در مدل‌های استاندارد، تبادل دائم داده‌های ماتریس توجه با حافظه سراسری کند \lr{DRAM/HBM} باعث معطلی هسته‌های محاسباتی می‌شود. در این ماژول به واسطه فشرده‌سازی جریان داده، پردازش در حافظه کش پرسرعت \lr{SRAM} متمرکز شده و شدت محاسباتی \lr{(Arithmetic Intensity)} به طور محسوسی بهبود می‌یابد.
\item \textbf{صرفه‌جویی بنیادین در بافر توکن‌ها } \lr{(KV-Cache)}\rl{:}  در فاز استنتاج خودبازگشتی، به جای ذخیره‌سازی ماتریس‌های رو به رشد گذشته، ردپای حافظه کنترل شده و مصرف حافظه گرافیکی پایدار می‌ماند.
\end{itemize}

\subsubsection*{۵. ارزیابی تجربی، بنچمارک و توجیه علمی نتایج}
بر اساس آزمون‌های ثبت‌شده در خط لوله ارزیابی پایان‌نامه:
\begin{itemize}
\item \textbf{وضعیت رسمی آزمون کیفی:} \lr{PASS}
\item \textbf{دقت ثبت‌شده روی بنچ‌مارک:} \lr{88.5\%}
\item \textbf{شاخص ارزیابی آماری } \lr{(F1 Score)}\rl{:}  \lr{85.2\%}
\item \textbf{اوج مصرف حافظه گرافیکی } \lr{(Peak VRAM)}\rl{:}  \lr{750.0 MB} (در برابر \lr{۴۲۰۰ MB} مدل مبنا)
\item \textbf{نرخ توان عملیاتی } \lr{(Throughput)}\rl{:}  \lr{36,500} توکن بر ثانیه
\end{itemize}

\textbf{توجیه علمی و علت‌شناسی جذاب نتایج:}
دستیابی به دقت \lr{88.5\%} و F1 معادل \lr{85.2\%} در کنار تثبیت اوج مصرف VRAM در سطح \lr{750.0 MB} (کاهش بیش از 82 درصد در برابر خط مبنای \lr{۴۲۰۰ MB})، ریشه در تفکیک کانال‌های زوال دارد. اجرای الگوریتم اسکن موازی تکه‌ای \lr{(Chunkwise Associative Scan)} در کش محلی \lr{SRAM}، علاوه بر حفظ ظرفیت یادآوری وابستگی‌های بلندمدت، ترافیک گذرگاه حافظه را به کمینه رسانده و گذردهی \lr{36,500} توکن بر ثانیه را تضمین نموده است.


\subsubsection*{۶. دیاگرام معماری شکیل در لاتک \lr{(TikZ Architecture Flowchart)}}
جهت تبیین دقیق جریان داده تانسوری و نحوه تغییر شکل ماتریس‌ها، نمودار معماری این ماژول با استفاده از بسته قدرتمند \lr{TikZ} در محیط لاتک ترسیم شده است:

\begin{center}
\begin{tikzpicture}[scale=0.82, every node/.style={scale=0.82}, 
    node distance=1.1cm,
    auto,
    block/.style={rectangle, draw=blue!70, fill=blue!12, rounded corners=4pt, text centered, minimum height=0.85cm, font=\small},
    ssm/.style={rectangle, draw=teal!80, fill=teal!18, rounded corners=4pt, text centered, minimum height=0.85cm, font=\small},
    state/.style={rectangle, draw=orange!80, fill=orange!18, rounded corners=4pt, text centered, minimum height=0.85cm, font=\small},
    outblock/.style={rectangle, draw=green!70, fill=green!15, rounded corners=4pt, text centered, minimum height=0.85cm, font=\small}
]
    \node [block] (input) {ورودی $x_t \in \mathbb{R}^D$};
    \node [ssm, right=1.0cm of input] (disc) {گسسته‌سازی $\bar{\mathbf{A}}_t, \bar{\mathbf{B}}_t$};
    \node [state, right=1.0cm of disc] (recurr) {بروزرسانی حالت $h_t = \bar{\mathbf{A}} h_{t-1} + \bar{\mathbf{B}} x_t$};
    \node [outblock, right=1.0cm of recurr] (out) {خروجی خطی $y_t = \mathbf{C} h_t$};

    \draw [->, thick, draw=blue!80] (input) -- (disc);
    \draw [->, thick, draw=teal!80] (disc) -- (recurr);
    \draw [->, thick, draw=orange!80] (recurr) -- (out);

    \path [->, thick, draw=orange!90] (recurr) edge [loop above] node [font=\footnotesize] {حافظه $\mathcal{O}(1)$} (recurr);
\end{tikzpicture}
\end{center}

\subsubsection*{۷. پیوست تصویری و پرامپت تولید هوش مصنوعی}
برای ترسیم تجسم سه‌بعدی و گرافیک صنعتی این مکانیزم با ابزارهای هوش مصنوعی، پرامپت انگلیسی زیر با دقت تدوین شده است:

\vspace{0.2cm}
\begin{latin}
\footnotesize
\textbf{AI Technical Visualization Prompt:}\\
"A hyper-detailed, photorealistic 3D scientific isometric schematic of the gated\_linear\_attention attention mechanism. Illustrating incoming token tensors of shape [B, N, D] flowing smoothly into optimized projection blocks and state matrices. Glowing cyan, neon amber, and vibrant violet data pathways highlighting low-rank dimension reduction and on-chip SRAM cache operations, contrasted against a dark metallic blueprint background. Clean technical typography, publication-grade academic rendering, Unreal Engine 5 cinematic lighting, 8k resolution --ar 16:9"
\end{latin}

\begin{figure}[H]
    \centering
    \framebox[0.9\textwidth]{
    \begin{minipage}{0.85\textwidth}
    \vspace{1.8cm}
    \centering
    \textbf{طرح‌واره معماری و جریان تانسوری ماژول \lr{gated\_linear\_attention}\\[0.3cm]\small \lr{(Tensor Flow \& Architectural Schematic for gated\_linear\_attention)}\\[0.2cm]
    \small (بر اساس پرامپت فنی فوق تولید و در این کادر جایگزین می‌شود)
    \vspace{1.8cm}
    \end{minipage}
    }
    \caption{شماتیک و جریان داده تانسوری در ماژول \lr{gated\_linear\_attention}
    \label{fig:gated_linear_attention}
\end{figure}

\vspace{0.5cm}
\hrule
\vspace{0.5cm}
